\section{Large fundamental groups}\label{section:Large.fundamental.groups}
In this section, we begin the study of fundamental groups of varieties that develop to the model of a Cartan geometry.
We build up this theory until we find that if the fibers of a fibration in a~Cartan geometry have solvable holonomy for the Cartan connection, then they are tori.

A connected complex projective variety $M$ has \emph{large fundamental group} if every connected positive-dimensional closed
complex subvariety $Z \subset M$ with normalization $Z_0 \to Z$ has infinite fundamental group image $\fundamentalGroup{Z_0} \to \fundamentalGroup{M}$ \cite{Kollar:1993}.

\begin{Lemma}\label{lemma:dev.large}
Suppose that $M$ is a smooth projective variety admitting a minimal geometry modelled on a complex homogeneous space $(X, G)$.
Suppose that \rationalMap{Z}{M} is a meromorphic map from a reduced connected compact complex space, and that this map develops into $X$.
Then the map extends to a holomorphic map $Z \to M$ lying in a leaf of the ant foliation on $M$ and the developing map takes $Z$ to a fiber of the ant fibration.
Suppose that $Z \to M$ is a finite map.
Then $Z$ has large fundamental group.
\end{Lemma}

\begin{proof}
Pullback the bundles $E|_Z = G|_Z$, and so identify the tangent bundles $ TX|_Z= TM|_Z$.
Since the model $X$ is homogeneous, its tangent bundle is spanned by global sections, and so its anticanonical bundle as well, and so nef.
The canonical bundle of $M$ is nef.
Therefore, the canonical bundle of $\cb{M}$ restricts to a nef line bundle on $Z$ with nef dual bundle the restriction of $\acb{X}$.
Hence $\cb{M}$ restricts to a trivial bundle on $Z$ as does $\cb{X}$.
Wedging sections, $ TX|_Z$ is spanned by global sections, but global sections of $ TX|_Z= TM|_Z$ can not vanish anywhere, so~${TX|_Z= TM|_Z}$ is a trivial bundle.
If a section of $TX$ vanishes at some point of $X$, then by homogeneity some such section vanishes at any point of $X$, in particular at some point of $Z$.
Any section of $TX$ which vanishes at some point of $Z$ vanishes at all points of $Z$, since $\acb{X}$ is trivial on $Z$.
Write the map $Z \to M$ as $f$, and its development $Z \to X$ as $\widehat{f}$.
So at every point of $\widehat{f}(Z)$, precisely the same vector fields on $X$ vanish and precisely the same sections of~$\acb{X}$.
In other words, $Z \to X$ lies in an anticanonical fiber.
All holomorphic functions on that fiber pull back to constants on $Z$, so $Z \to X$ lies in an ant fiber.
So $Z \to M$ lies in an ant leaf.
Each ant fiber $(X_0, G_0)$ is a connected complex homogeneous space, and its universal covering space is a Stein complex Lie group by
Lemma~\ref{lemma:ant.Stein}.
So $Z \to X$ lifts to a holomorphic map to that Stein complex Lie group from some covering space of $Z$.

Suppose that $g \colon Y \to Z$ is a holomorphic map from a connected compact normal complex space.
Under the group morphism $\fundamentalGroup{Y} \to \fundamentalGroup{Z}$, suppose further that the preimage of~$\fundamentalGroup{\widehat{Z}}$ is a finite index subgroup of $\fundamentalGroup{Y}$.
After replacing $Y$ by a finite unramified covering space, we can assume that this image is trivial.
The map $Y \to Z \to X_0$ lifts to $Y \to \widehat{X}_0$, which is a~Stein manifold.
Hence $Y \to \widehat{X}_0$ is constant and so $Y \to X_0$ and $Y \to M$ are constant so~$Y$ is a point.
In particular, if $\fundamentalGroup{Y} \to \fundamentalGroup{Z}$ is finite, $Y$ is a point.
\end{proof}

\begin{Example}
A non-uniruled smooth projective variety with a holomorphic parabolic geometry has trivial ant foliation, since its model is a flag variety, hence has trivial ant fibration.
Therefore, a meromorphic map from a reduced connected compact complex space $Z$ develops to a minimal geometry if and only if that map is constant.
So it develops to an arbitrary holomorphic parabolic geometry, on a connected manifold, just when it lies in a fiber of the fibration to the associated minimal geometry.
Roughly speaking, once we drop out the rational curves, we can not develop anything.
\end{Example}

\begin{Lemma}\label{lemma:large}
Suppose that $M$ is a smooth projective variety bearing a minimal geometry modelled on a complex homogeneous space $(X, G)$.
Suppose that \rationalMap{Z}{M} is a meromorphic map from a positive-dimensional reduced connected compact complex space.
Then this meromorphic map extends uniquely to a holomorphic map.
Suppose that $Z \to M$ is a finite map.
Take the normalization $Z_0 \to Z \to M$.
Either
\begin{enumerate}\itemsep=0pt
\item[$(1)$]
the image of $\fundamentalGroup{Z_0} \to \fundamentalGroup{M}$ is infinite or
\item[$(2)$]
$Z$ has large fundamental group.
\end{enumerate}
\end{Lemma}

\begin{proof}
Meromorphic maps to $M$ from reduced compact complex spaces extend to holomorphic maps \cite[Lemma 26]{Biswas.McKay:2016}.
Replace $Z$ by its normalization.

Suppose that the image of $\fundamentalGroup{Z} \to \fundamentalGroup{M}$ is finite.
After replacing $M$ by a finite covering space, we can arrange that the image of $\fundamentalGroup{Z}
\to \fundamentalGroup{M}$ is trivial.

Suppose that $Z \to M$ is finite.
By Remmert's proper mapping theorem \cite[p.~213]{Grauert.Remmert:1984}, the image of $Z
 \to M$ is a projective subvariety;
pick a proper curve inside the image of~${Z \to M}$.
Replace $Z$ by the preimage of that curve inside $Z$.
The curvature of the Cartan geometry ${E \to M}$ vanishes on $Z$ because the curvature is a holomorphic $2$-form, so vanishes on curves.
Since $\fundamentalGroup{Z} \to \fundamentalGroup{M}$ has trivial image, the connection has no holonomy.
The Cartan connection integrates to a developing map to the model $Z \to X$.
By Lemma~\ref{lemma:dev.large}, the curve $Z$ has large fundamental group.
Now if $Z$ is not a curve, and contains a subvariety whose normalization has
finite fundamental group image in $Z$, pick a curve in that subvariety.
\end{proof}

\begin{Corollary}\label{corollary:tori}
Suppose that $M$ is a smooth projective variety bearing a minimal geometry.
Take a smooth projective variety $Z$ with trivial canonical bundle.
If there is a finite meromorphic map \rationalMap{Z}{M} then $Z$ has a finite unramified covering by an abelian variety.
\end{Corollary}

\begin{proof}
Meromorphic maps to $M$ from reduced compact complex spaces extend to holomorphic maps \cite[Lemma 26]{Biswas.McKay:2016}.
After perhaps replacing by a finite unramified covering, any smooth projective variety $Z$ with trivial canonical bundle splits into a product $Z=A \times B$ where $A$ is an abelian variety and $B$ is a simply connected Calabi--Yau manifold \cite{Bogomolov:1974}.
Replace $Z$ by $\{a_0\} \times B$ for any $a_0 \in A$ to arrange that $Z$ is simply connected, contradicting
Lemma~\ref{lemma:large} unless $B$ is a~point.
\end{proof}

\begin{Lemma}\label{lemma:solvable.nef}
Suppose that $M$ is a connected compact K\"ahler manifold bearing a minimal geometry with solvable holonomy and either $M$ has nef canonical bundle or $M$ is Moishezon.
Then $M$ is a complex torus, after perhaps replacing $M$ by a finite unramified covering.
\end{Lemma}

\begin{proof}
By Lie's theorem (see \cite[Theorem 3]{Serre:2001}), every associated vector bundle of the bundle $E_G \to M$ has a filtration into a flag of holomorphic vector bundles, with quotient line bundles.
Those line bundles admit holomorphic connections, so are nef.
In particular, the adjoint bundle~$E_G \times^G \LieG$ has such a filtration.
Therefore, the adjoint bundle is an extension of nef line bundles, and so is a nef vector bundle.
Since $TM$ is a quotient bundle of the adjoint bundle:
\[
E_G \times^G \LieG = E \times^H \LieG \to E \times^H \pr{\LieG/\LieH} = TM,
\]
$\acb{M}$ is nef.
But $\cb{M}$ is nef (by Theorem~\ref{theorem:drop} if $M$ is Moishezon), so $\cb{M} = 0$, after perhaps replacing by a finite unramified covering, hence a complex torus after perhaps another finite unramified covering \cite{Biswas.McKay:2010}.
\end{proof}

Take a complex homogeneous space \pr{X,G} and a point $x_0 \in X$.
Take the ant fibration $X = G/H \to \overline{X} = G/\overline{H}$, with fiber \pr{X_0,\GZ} where
\[
X_0 = \overline{H}/H = \GZ/\Gamma
\]
with $\GZ = \overline{H}/K$ and $\Gamma = H/K$ where
\[
K = \bigcap_{g\in\overline{H}} gHg^{-1}
\]
contains $H^0$.
Denote the universal covering group as $\tGZ \to \GZ$ and let $\widetilde{\Gamma} \subset \tGZ$ be the
preimage of $\Gamma \subset \GZ$.

Take a complex Lie algebra morphism
$
\widetilde{f} \colon \C^{n} \to \LieGZ
$
from the abelian Lie algebra $\C^{n}$ and exponentiate into a complex Lie group morphism
$
\widetilde{f} \colon \C^{n} \to \GZ.
$
Suppose that $\Lambda \subset \C^{n}$ is a~lattice and that $\widetilde{f}\Lambda \subset \widetilde{\Gamma}$.
Let $T := \C^{n}/\Lambda$.
Quotient to a map $f \colon T \to X_0$.
Take an element~${g \in \GZ}$.
The map $gf \colon T \to X$ is a \emph{hoop}.

\begin{Lemma}\label{lemma:is.a.hoop}
Every holomorphic map from a connected compact K\"ahler manifold with trivial canonical bundle to a complex homogeneous space which pulls back the canonical bundle to be trivial is a hoop.
\end{Lemma}

\begin{proof}
Take a connected compact K\"ahler manifold $Z$ with trivial canonical bundle and a holomorphic map $f \colon Z \to X$ so that $f^* \cb{X}=0$.
Clearly $f(Z)$ lies in an anticanonical fiber of~$X$, so without loss of generality we can replace $X$ by that fiber, and so assume that $X=G/\Gamma$ where $\Gamma \subset G$ is a discrete subgroup.
Since $Z$ is compact, every holomorphic function on $X$ pulls back to a constant on $Z$, so $Z$ lies in an ant fiber of $X$.
Again, without loss of generality, replace $X$ by that fiber.
After perhaps replacing $Z$ by a finite unramified covering space, split~$Z = T \times B$ where $T$ is a complex torus and $B$ is a simply connected compact K\"ahler manifold with trivial canonical bundle.
For each $t_0 \in T$ the map $f$ lifts to take $\{t_0\} \times B$ to~$\widetilde{X}_0$, a simply connected complex Lie group, so $B$ is a point.
The map $f \colon T \to X$ pulls back the right invariant Maurer--Cartan form ${\rm d}g g^{-1}$ to be $f^* {\rm d}g g^{-1}
 = c {\rm d}z$, a constant multiple of the translation invariant coframing.
The holomorphic vector fields on $T$, all of which are constant, map by $f'$ to holomorphic sections of the pullback tangent bundle (also constant), representing tangent vector fields, taking brackets (all zero) to brackets, so $f$ is a morphism of Lie algebras from an abelian Lie algebra.
Therefore, the lift $\widetilde{f} \colon \C^{n} \to G$ which satisfies $\widetilde{f}(0)=1$ (after suitable translation) is a complex Lie group morphism.
\end{proof}
\begin{Corollary}\label{corollary:hoops}
Suppose that $M$ is a non-uniruled smooth projective variety bearing a holomorphic Cartan geometry.
Take a finite map $Z \to M$ which develops to $X$, from a smooth projective variety $Z$ with trivial canonical bundle.
Then, up to replacing $Z$ by a finite unramified covering map, $Z$ is an abelian variety, and the developing map is a hoop.
\end{Corollary}

\begin{Example}
If $M = \C^{n}/\Lambda_0$ is a complex torus and $X = \C^{n}/\Lambda_1$ is another complex torus and $G = X$, then $M$ has the obvious flat \pr{X,G}-geometry, unique up to linear transformation, with identity map as developing map.
If, after such a linear transformation, we can find a linear subspace $V \subset \C^{n}$ containing a
lattice $\Lambda \subset \Lambda_0 \cap \Lambda_1$, then $V/\Lambda$ develops to the model by the identity map.
\end{Example}

\begin{Lemma}\label{lemma:abelian.pancake}
Take a non-uniruled smooth projective variety $M$ with a holomorphic Cartan geometry $E \to M$ with model $(X, G)$.
Take a rational dominant map $\rationalMap{M}{S}$.
Suppose that on the general fiber of that map, the pullback bundle of $E_G$ admits a holomorphic connection, perhaps not equal to the Cartan connection, with holonomy acting on $\LieG$ as a solvable group of linear transformations.
Then the general fiber admits a finite unramified covering by an abelian variety.
Moreover, $\cb{M}$ is trivial on the general fiber.
If the connection is flat, then the holonomy on each fiber is abelian.
\end{Lemma}

\begin{proof}
Take a general fiber $M_{s} $ of $\rationalMap{M}{S}$ and let $E_{G_{\rm s} } := E_G|_{M_{s} }$.
By Lie's theorem \cite[Theorem 3]{Serre:2001}, the adjoint vector bundle $\ad{E_{G_{\rm s} }} := \amal{E_{G_{\rm s}}}{G}{\LieG}$ of the bundle $E_{G_{\rm s} }
\to M_{s} $ has a filtration into a flag of holomorphic vector bundles, with quotients being line bundles.
The holomorphic connection induces a holomorphic connection on each quotient line bundle, so each quotient line bundle is nef.
Therefore, the adjoint bundle is an extension of nef line bundles, and so is a nef vector bundle.
Since the tangent bundle $TM|_{M_{s} }$ is a quotient bundle of the adjoint bundle:
\[
\amal{E_G}{G}{\LieG} = \amal{E}{H}{\LieG} \to \amal{E}{H}{\pr{\LieG/\LieH}} = TM,
\]
$\acb{M}$ is nef along $M_{s} $.
But $\cb{M}$ is nef on $M$, since $M$ is not uniruled, so also nef on $M_{s} $ and so $\cb{M} = 0$ on $M_{s} $.

The general fiber $M_{s} $ is smooth.
By adjunction, the anticanonical bundle $\acb{M_{s}} = - \cb{M} |_{M_{s}}$ is nef.
But $\cb{M_{s}}$ is also nef as every subvariety of $M$ is minimal by Theorem~\ref{theorem:drop}.
Hence the general fiber $M_{s}$ has trivial canonical bundle.
By Corollary~\ref{corollary:tori}, $M_{s}$ admits a finite unramified covering by an abelian variety, i.e., the general fiber of $\rationalMap{M}{S}$ admits a finite unramified covering by an abelian variety.
Since the fundamental group of any abelian variety is abelian, its image under the holonomy morphism is also abelian.
\end{proof}

\section{Review of the Shafarevich fibration}\label{section:Shafarevich}
In this section, we review the definition and basic properties of the Shafarevich fibration (a well known construction in algebraic geometry).

Take any compact K\"ahler manifold $M$ equipped with a morphism $\rho \colon \fundamentalGroup{M} \to G$ of groups where $G$ is a complex semisimple Lie group.
A \emph{Shafarevich map} for $\rho$ is a meromorphic map~${\Shaf[\rho]{} \colon \rationalMap{M}{\Shaf[\rho]{M}}}$ to a~normal complex space, surjective and with connected fibers, so that a~connected closed subvariety $Z \subset M$ through a very general point is mapped to a~point by the Shafarevich map just when the normalization of $Z$ has fundamental group with finite image inside $G$.
Every group morphism $\rho \colon \fundamentalGroup{M} \to G$ with Zariski dense image in a~complex semisimple Lie group $G$ has a Shafarevich map \cite[Definition~2.13]{Campana/Claudon/Eyssidieux:2015}, \cite[p.~105]{Zuo:1999} with the additional property that, after replacing $M$ by a finite covering space, the target $\Shaf[\rho]{M}$ is bimeromorphic to the total space of a smooth fibration of complex tori over an algebraic variety of general type \cite[Theorem~4]{Campana/Claudon/Eyssidieux:2015}.

Suppose that $\Sigma \subset \fundamentalGroup{M}$ is the kernel of a representation of $\fundamentalGroup{M}$.
Also after replacing $M$ by a finite covering space, there is a map, also called a Shafarevich map,
$
\Shaf[\Sigma]{} \colon \rationalMap{M}{\Shaf[\Sigma]{M}}
$
which contracts to a point precisely those subvarieties $Z \subset M$ passing through a very general point whose fundamental group has a finite index subgroup lying inside the given subgroup $\Sigma \subset \fundamentalGroup{M}$ \cite[Theorem 1.3]{Katzarkov:1999}, \cite{Katzarkov:1997}; for compact K\"ahler manifolds the existence of this Shafarevich map follows
from \cite[Theorem 4]{Campana/Claudon/Eyssidieux:2015} and \cite{Katzarkov:1997}.
A representation $\rho \colon \fundamentalGroup{M} \to G$ of the fundamental group of a compact K\"ahler manifold is \emph{big} if the associated Shafarevich map is a birational map.

Take a complex Lie group $G$, with identity component $G^0$, whose component group $G/G^0$ admits a faithful finite-dimensional representation.
Then $G$ admits a faithful finite-dimensional holomorphic representation just when $G^0$ is the semidirect product of a simply connected solvable Lie group and a connected reductive complex
Lie group \cite[Theorem~16.3.7]{Hilgert.Neeb:2012}.

Take a compact irreducible reduced complex space $M$ containing no rational curves.
Suppose that $X$ is a compact irreducible reduced complex space. Then every meromorphic map $\rationalMap{X}{M}$ extends to a unique holomorphic map $X \to M$, by Hironaka's Chow lemma \cite[Corollary~2.9]{Grauert:1994}.

The \emph{Iitaka conjecture}, also called the \emph{$C_{mn}$ conjecture}, for a dominant rational morphism $\rationalMap{X}{Y}$ of projective varieties with connected fibers, claims that $\kappa_X \ge \kappa_Y + \kappa_{X_y}$ for the general fiber $X_y$.
The Iitaka conjecture is verified for surjective morphisms for which all sufficiently high powers of the canonical bundle of the general fiber are spanned by global sections \cite[Corollary~1.2]{Kawamata:1985}.

